\section{Course Overview and Modeling Strategy}

Higher brain function is a modeling problem before it is a tooling problem. Memory, planning, inference, abstraction, and decision making are all latent processes: we rarely observe them directly, but instead infer them from behavior, neural activity, and task structure. The course therefore asks how to build models whose hidden variables are meaningful, constrained, and experimentally testable.

The working standard is that every model should answer four questions. What is the observable behavior? Which latent variables explain it? What update rule moves those variables through time? What would count as evidence against the model? Keeping these questions explicit prevents the notes from becoming a catalog of architectures without interpretation.

\subsection{What makes a model useful?}

A model of higher brain function should do more than fit observations. It should compress behavior into a tractable set of assumptions, explain why certain error patterns appear, and suggest new experiments or interventions. In cognitive and systems neuroscience, this usually means balancing three goals:
\begin{itemize}
  \item computational clarity,
  \item biological plausibility,
  \item empirical testability.
\end{itemize}

\begin{figure}[h]
\centering
\includegraphics[width=0.82\textwidth]{figures/mhbf-model-triangle.png}
\caption{Why useful models live in tension. Strong higher-brain models stay interpretable, stay anchored to plausible mechanisms, and still make contact with measurable data.}
\end{figure}

\paragraph{Practical heuristic.}
If a model explains everything equally well, it usually explains nothing. Useful models make sharp commitments about mechanisms, variables, and expected failures.

\subsection{A common abstraction stack}

Many courses on higher brain function move across a recurring abstraction ladder: behavior and task structure, latent cognitive variables, circuit or algorithmic mechanisms, and normative or optimization principles.

That ladder connects classical cognitive models to modern neural-network and dynamical-systems viewpoints. Reinforcement learning, attractor dynamics, and recurrent computation all make the same point in different mathematical languages: representation, dynamics, and objective functions have to be studied together.

\begin{figure}[h]
\centering
\includegraphics[width=0.82\textwidth]{figures/mhbf-abstraction-stack.png}
\caption{The abstraction ladder used throughout the course. Explanations move downward from observed behavior to hidden variables and mechanism, then upward again to ask whether the mechanism reproduces the behavior.}
\end{figure}

\paragraph{Generic state-space view.}
\[
x_{t+1} = f(x_t, u_t; \theta)
\]

This form is broad enough to include recurrent neural networks, drift-diffusion style updates, working-memory attractors, and many cognitive-state models.

\begin{figure}[h]
\centering
\includegraphics[width=0.82\textwidth]{figures/mhbf-state-space.png}
\caption{State-space intuition. The hidden state carries forward what the system currently ``knows,'' inputs perturb that state, and the update rule determines the next internal configuration.}
\end{figure}

\subsection{Reusable model code}

Executable examples should make modeling assumptions visible. Useful code modules include simulation utilities for canonical cognitive tasks, fitting routines for simple latent-state models, plotting helpers for trajectories and phase portraits, and small reading replications that test whether a model reproduces the qualitative phenomenon it claims to explain.

\subsection{Next steps}

Each conceptual block should end with a compact model checklist: the state variables, the objective or update rule, the measurable predictions, and the most likely confound. This makes later comparisons between attention, memory, decision, and control models precise rather than purely verbal.
